\documentclass[12pt]{article}

\usepackage{sbc-template}
\usepackage{graphicx,url}
\usepackage[utf8]{inputenc}
\usepackage[brazil]{babel}

\usepackage{graphicx}
\usepackage{svg}
\usepackage{amsmath}
\usepackage{amssymb}
\usepackage{booktabs}
\usepackage{array}
\usepackage{rotating}
\usepackage{tabularx}
\usepackage{float}
\usepackage{soul}
\usepackage[section]{placeins}
\usepackage{longtable}
%\usepackage[latin1]{inputenc}  

\graphicspath{{Images/}}

\sloppy

\title{Arquitetura Multicamadas para Segurança em Modelos de Linguagem de Grande Escala (LLMs)}

\author{Juan G. L. Cavalcante\inst{1}, Lucas E. S. Barros\inst{1}, João Vitor\inst{1}, Joás Silva\inst{1}, Lucas Messias\inst{1} }


\address{Bacharel em Ciências da Computação -- Universidade Federal do Agreste \\
    de Pernambuco (UFAPE) \\
   55292-270 -- Garanhuns -- PE -- Brasil
}

\begin{document} 

\maketitle

\begin{abstract}
This paper presents a multilayered architecture for security in Large Language Models (LLMs), integrating sanitization, threat detection, bias mitigation, and post-generation validation. The conducted experiments demonstrate high effectiveness against explicit attacks and persistent challenges in detecting linguistic biases, contributing to advancements in AI safety and governance.

\end{abstract}
     
\begin{resumo} 
  Este artigo apresenta uma arquitetura multicamadas para segurança em Modelos de Linguagem de Grande Escala (LLMs), integrando sanitização, detecção de ameaças, mitigação de vieses e validação pós-geração. Os experimentos conduzidos demonstram alta eficácia contra ataques explícitos e desafios persistentes na detecção de vieses linguísticos, contribuindo para avanços em segurança e governança de IA.
\end{resumo}


\section{Introdução}
\label{sec:introduction}

Modelos de Linguagem de Grande Escala (LLMs) têm se tornado componentes centrais em aplicações modernas, incluindo assistentes virtuais, sistemas conversacionais, plataformas educacionais e ambientes corporativos. Sua capacidade de gerar, interpretar e transformar linguagem natural em larga escala permite níveis de automação antes inviáveis. Entretanto, essa mesma flexibilidade introduz riscos significativos relacionados à segurança, ética e confiabilidade, especialmente quando tais modelos são expostos diretamente a usuários finais ou integrados a sistemas críticos.

Diversas vulnerabilidades inerentes aos LLMs têm sido amplamente documentadas na literatura, incluindo ataques de \textit{prompt injection}, tentativas de \textit{jailbreak} para subverter instruções do sistema, geração de conteúdo malicioso, vieses socioculturais reproduzidos ou amplificados pelo modelo e alucinações (\textit{hallucinations}) que podem resultar em informações factualmente incorretas ou perigosas. Além disso, LLMs podem inadvertidamente vazar informações pessoais identificáveis (PII), expondo organizações e usuários a riscos legais e de privacidade.

Soluções tradicionais baseadas apenas em filtragem estática, listas de bloqueio ou simples verificações sintáticas têm se mostrado insuficientes para lidar com ameaças cada vez mais sofisticadas. Isso ocorre porque muitos ataques exploram dependências semânticas, ambiguidades linguísticas e lacunas nos mecanismos de alinhamento do modelo, exigindo abordagens mais abrangentes e sensíveis ao contexto.

Diante desse cenário, este trabalho propõe uma arquitetura multicamadas para segurança de LLMs, composta por módulos de sanitização estrutural, detecção de ameaças explícitas, identificação de vieses e validação de saídas, orquestrados de forma sistemática por um serviço centralizado. O objetivo é oferecer um mecanismo de defesa capaz de atuar em diferentes níveis do fluxo de interação, mitigando riscos antes e depois da geração da resposta pelo modelo.

Para avaliar a eficácia da solução proposta, foram conduzidos experimentos envolvendo 30 casos de teste distribuídos entre categorias críticas, incluindo ataques maliciosos, tentativas de \textit{jailbreak}, exposição de PII e múltiplas subcategorias de viés linguístico. Métricas quantitativas — como precisão, recall, acurácia e F1-score — foram utilizadas para mensurar o desempenho da arquitetura em cenários adversariais. Os resultados indicam desempenho superior nas categorias de risco explícito, enquanto apontam desafios persistentes na detecção de vieses implícitos.

As principais contribuições deste trabalho são:

\begin{itemize}
    \item a proposição de uma arquitetura multicamadas para segurança em LLMs, combinando verificações sintáticas, estruturais e semânticas;
    \item o desenvolvimento de módulos especializados para sanitização, bloqueio de ataques, detecção de vieses e validação pós-geração;
    \item a definição de um fluxo padronizado de comunicação que aumenta a interpretabilidade e auditabilidade das decisões de segurança;
    \item a avaliação experimental da arquitetura em um conjunto diversificado de cenários adversariais, com análise quantitativa e comportamental dos resultados.
\end{itemize}

O restante deste artigo está estruturado da seguinte forma. A Seção~\ref{sec:background} apresenta os fundamentos teóricos relacionados a LLMs e mecanismos de segurança. A Seção~\ref{sec:related} discute os principais trabalhos relacionados. A
Seção~\ref{sec:architecture} descreve a arquitetura do sistema. A Seção~\ref{sec:proposed_model} apresenta o modelo conceitual subjacente. A Seção~\ref{sec:resultados_finais} discute os resultados experimentais. Por fim, a Seção~\ref{sec:conclusion} apresenta as conclusões e trabalhos futuros.



\section{Trabalhos Relacionados}
\label{sec:related}

Para a construção do referencial teórico deste trabalho, realizou-se uma análise sistemática de publicações recentes que investigam a segurança de prompts em Modelos de Linguagem de Grande Escala (LLMs). Os estudos examinados abrangem propostas de frameworks de defesa, sistematizações de conhecimento, mecanismos baseados em consultas estruturadas, análises empíricas de vulnerabilidades e avaliações voltadas à privacidade e percepção dos usuários. Essa revisão permitiu delinear o estado da arte, evidenciando o progresso alcançado na compreensão e mitigação de ataques de injeção de prompt e ameaças correlatas.

Entretanto, a análise comparativa revelou que cada um desses trabalhos apresenta limitações relevantes. Como pode ser observado na Tabela 1:

\small
\setlength{\extrarowheight}{2pt}

\begin{longtable}{
      >{\raggedright\arraybackslash}p{0.22\textwidth}
      >{\raggedright\arraybackslash}p{0.18\textwidth}
      >{\raggedright\arraybackslash}p{0.18\textwidth}
      >{\raggedright\arraybackslash}p{0.18\textwidth}
      >{\raggedright\arraybackslash}p{0.20\textwidth}}
\caption{Comparison of Related Works on Prompt Injection Defense in LLMs}
\label{tab:related_works}
\\
\toprule
\textbf{Referencia} & \textbf{Foco} & \textbf{Metodologia} & \textbf{Métricas} & \textbf{Lacuna} \\
\midrule
\endfirsthead

      Lin, Huawei \emph{et al.} \cite{Lin2025UniGuardian} 
      & Proposta de defesa unificada (UniGuardian) contra ataques de LLM (injeção, backdoor, adversariais).
      & Propõe um framework de detecção unificado para analisar prompts e saídas.
      & (1) auROC (Area Under the Receiver Operator Characteristic Curve) e (2) auPRC (Area Under the Precision-Recall Curve).
      & Falta de profundidade no que tange a ataques de backdoor mais complexos ou ofuscados, gerando falsos positivos. \\
    \midrule
      Hong, Hanbin \emph{et al.} \cite{Hong2025SystematizationOfKnowledge} 
      & Sistematização do conhecimento em segurança de prompts em LLMs, propondo uma taxonomia e um conjunto de métricas próprias para padronizar avaliações de ataques e defesas.
      & Revisão sistemática e proposta de uma taxonomia multi-nível. Desenvolvimento de toolkit aberto e dataset (JailbreakDB) para avaliação padronizada.
      & Propõe métricas padronizadas próprias integradas a um toolkit de avaliação com taxonomia hierárquica e perfis de ameaça formais.
      & Limitação na aplicação prática das métricas e taxonomias propostas; necessidade de validação contínua e ampliação para modelos multimodais e cenários reais. \\
    \midrule
      Chen, Sizhe \emph{et al.} \cite{Chen2025StruQ}
      & Proposta de defesa chamada \textit{StruQ}, que utiliza consultas estruturadas (\textit{structured queries}) para separar dados do prompt, reduzindo o risco de injeções maliciosas em LLMs.
      & Implementação experimental com análise de desempenho e eficácia em diferentes cenários de injeção. Avalia a separação entre dados e instruções para mitigar ataques.
      & Taxa de sucesso de ataque, latência e sobrecarga do sistema (\textit{defense overhead}).
      & Boa eficácia em prompts simples, mas limitação em ataques indiretos e cenários complexos de múltiplas etapas, exigindo integração com outras técnicas de defesa. \\
      
    \midrule
      Benjamin, Victoria \emph{et al.} \cite{benjamin2024}
      & Análise sistemática da vulnerabilidade de 36 LLMs a ataques de injeção de prompt focados em gerar código de keylogger.
      & 4 prompts de injeção direta contra 36 LLMs. Análise estatística (Correlação, Random Forest, SHAP, PCA).
      & (1) Taxa de sucesso (2) Importância de features. (3) Correlação entre prompts.
      & Necessidade de testes multilíngues, múltiplas etapas de complexidade. \\
    \midrule
      Sebastian, Glorin. \cite{Sebastian2023SecurityUserInformation:}
      & Investigar a proteção de dados e privacidade em chatbots (foco no ChatGPT). Avaliar Tecnologias de Aprimoramento de Privacidade (PETs).
      & Revisão da literatura, análise de técnicas (ex: privacidade diferencial) e uma pesquisa (survey) com 177 usuários.
      & Métricas de percepção da pesquisa: (1) Nível de preocupação, (2) Disposição para sacrificar desempenho, (3) Consciência sobre vazamentos.
      & Análise de riscos à privacidade, preocupação dos usuários através de pesquisas. \\
    \midrule
    \textbf{This Work}
      & \textbf{Análise comparativa de frameworks de defesa contra injeção de prompt em LLMs.}
      & \textbf{Prova de conceito e simulações baseadas em diferentes cenários de ataque.}
      & \textbf{Taxa de sucesso de ataque, resiliência e custo computacional.}
      & \textbf{Busca preencher a lacuna de falta de metodologia padronizada para comparação empírica entre defesas.} \\
    \bottomrule
\end{longtable}

Nesse sentido, as lacunas observadas justificam a necessidade de investigações adicionais que consolidem métricas, experimentos e metodologias, oferecendo uma base mais consistente para o desenvolvimento e comparação de soluções de segurança em LLMs.

\section{Metodologia}
\label{sec:methodology}

A metodologia adotada neste trabalho foi estruturada de forma a avaliar, de maneira sistemática e controlada, a eficácia da arquitetura proposta para segurança em Modelos de Linguagem de Grande Escala (LLMs). O processo metodológico integra a geração de cenários de ataque, a classificação de riscos, a execução de testes em múltiplos idiomas e a análise quantitativa dos resultados, permitindo uma investigação abrangente das capacidades e limitações do sistema.

\subsection{Levantamento do estado da arte}

Para fundamentar a investigação proposta, realizou-se um levantamento sistemático do referencial teórico em bases de dados amplamente reconhecidas na área de Computação e Inteligência Artificial, incluindo Google Acadêmico, IEEE Xplore e a Plataforma Lattes. A busca foi conduzida utilizando um conjunto de palavras-chave relacionadas ao tema central deste estudo, tais como “segurança de prompts”, “prompt injection”, “LLM security” e “prompt-based attacks”.

Inicialmente, os resultados recuperados foram triados com base nos títulos e resumos, de modo a identificar pesquisas diretamente alinhadas à temática de segurança em Modelos de Linguagem de Grande Escala (LLMs). Em seguida, procedeu-se à análise exploratória dos textos selecionados, considerando critérios como relevância metodológica, atualidade das abordagens e potencial contribuição para a compreensão do estado da arte.

Após esse processo de filtragem, cinco estudos foram selecionados para análise detalhada no âmbito deste projeto. Esses trabalhos representam diferentes perspectivas e metodologias aplicadas à mitigação de ataques de injeção de prompt e à avaliação de mecanismos de segurança em LLMs, fornecendo uma base sólida para a discussão crítica e para o desenvolvimento da abordagem proposta.

\subsection{Desenvolvimento da Prova de Conceito (POC)}

Após a elaboração do referencial teórico e a identificação das abordagens mais relevantes na área de segurança de prompts, procedeu-se ao desenvolvimento de um conjunto de aplicações destinadas à validação prática das técnicas analisadas. Esse desenvolvimento ocorreu de forma incremental, permitindo a implementação, teste e aprimoramento progressivo das funcionalidades relacionadas à detecção e mitigação de injeções de prompt em Modelos de Linguagem de Grande Escala (LLMs).

O ambiente experimental foi estruturado utilizando Python como linguagem base, devido à sua ampla adoção em projetos de inteligência artificial e ao ecossistema robusto de bibliotecas voltadas para processamento linguístico. Para a construção dos serviços de backend e exposição de interfaces de validação, empregou-se o framework FastAPI, que possibilita a criação de APIs de alto desempenho com tipagem clara e integração facilitada com modelos de IA.

Com o objetivo de garantir reprodutibilidade, isolamento e facilidade de implantação do ambiente, todas as aplicações foram containerizadas com Docker, permitindo a execução consistente dos serviços independentemente da infraestrutura subjacente. Além disso, integrou-se o framework Guardrails AI para implementar mecanismos de verificação, sanitização e conformidade dos prompts, assegurando que as entradas e saídas dos modelos fossem avaliadas segundo critérios de segurança previamente definidos.

\subsection{Construção dos Casos de Teste}

Para validar a assertividade da POC, foram desenvolvidos 30 casos de teste distribuídos entre quatro grupos principais:

\begin{itemize}
    \item \textbf{Ataques explícitos}: tentativas de \textit{prompt injection}, indução maliciosa e manipulação direta do comportamento do modelo;
    \item \textbf{Detecção de PII}: solicitações envolvendo CPF, e-mail e telefone;
    \item \textbf{Vieses linguísticos}: categorias relacionadas a gênero, raça, idade, orientação, localização geográfica;
    \item \textbf{Alucinações}: prompts projetados para induzir respostas factualmente inconsistentes.
\end{itemize}

\subsection{Sistematização dos Resultados e elaboração do Artigo}

Após a execução dos experimentos, os resultados obtidos foram organizados e analisados de forma sistemática, permitindo identificar padrões, comparar diferentes abordagens de defesa e avaliar a eficácia das estratégias de mitigação de ataques de prompt em LLMs. As métricas de desempenho, latência e assertividade foram tabuladas, destacando-se tanto os sucessos quanto as limitações observadas em cenários específicos.

Com base nessa análise, os achados foram estruturados de maneira lógica para compor o manuscrito, articulando a revisão de literatura, a metodologia aplicada e os resultados obtidos de forma coerente e integrada. A redação do artigo seguiu princípios de clareza, objetividade e rigor acadêmico, buscando apresentar as contribuições do estudo, as lacunas identificadas na literatura existente e as implicações práticas e teóricas das soluções desenvolvidas.


\input{sections/results}

\section{Conclusão e Trabalhos Futuros}
\label{sec:conclusion}

Este trabalho apresentou uma arquitetura multicamadas voltada para a segurança de Modelos de Linguagem de Grande Escala (LLMs), integrando mecanismos de sanitização, detecção de riscos explícitos, identificação de vieses e validação pós-geração. Os experimentos realizados demonstraram que a abordagem é tecnicamente viável, resiliente e capaz de mitigar de forma eficaz uma variedade de ameaças associadas ao uso de LLMs em ambientes abertos.

O sistema mostrou desempenho robusto na detecção de ataques explícitos, incluindo tentativas de \textit{prompt injection}, \textit{jailbreak} e solicitações maliciosas, além de apresentar eficácia total em todas as categorias de informações pessoais identificáveis (PII). As métricas obtidas — com F1-score igual a 1.00 para categorias de risco explícito — validam a capacidade da arquitetura de atuar como um mecanismo confiável de defesa em camadas. A ausência de falsos positivos reforça a minimização de interferência indevida na experiência do usuário.

Por outro lado, os resultados também evidenciaram limitações importantes, sobretudo na detecção de vieses linguísticos implícitos. Categorias como gênero, corpo, educação e LGBTQIA+ permaneceram com baixa ou nula taxa de sucesso, indicando que heurísticas baseadas apenas em regras e padrões simples não são suficientes para capturar nuances semânticas presentes em conteúdos enviesados. Essa limitação é consistente com desafios identificados na literatura e reforça a necessidade de incorporar mecanismos mais avançados de análise.

De forma geral, os resultados confirmam a hipótese central deste trabalho: \textit{uma arquitetura multicamadas, sensível ao contexto e baseada em verificações estruturais e semânticas, aumenta significativamente a segurança e a confiabilidade de sistemas baseados em LLMs}. O modelo proposto estabelece um arcabouço sólido para futuras extensões e aprimoramentos em sistemas de IA responsáveis.

\subsection*{Trabalhos Futuros}

Com base nas limitações observadas e nas oportunidades identificadas ao longo deste estudo, destacam-se as seguintes direções para trabalhos futuros:

\begin{itemize}
    \item \textbf{Aprimoramento da detecção de vieses}: Investigar o uso de classificadores supervisionados, embeddings semânticos e modelos específicos para vieses, a fim de ampliar a cobertura e reduzir falsos negativos.
    
    \item \textbf{Otimização do Bias Guardrail}: Reduzir a latência do módulo por meio de técnicas como pré-filtragem, paralelização, cache semântico e modelos mais leves para análise textual.

    \item \textbf{Ampliação do conjunto de testes}: Construir um dataset mais extenso, diversificado e representativo, incluindo casos reais de ataques, vieses complexos e cenários ambíguos.

    \item \textbf{Integração com modelos de avaliação contextual}: Incorporar módulos capazes de analisar contexto discursivo, pragmático e situacional, permitindo identificar violações implícitas que escapam a regras estáticas.

    \item \textbf{Melhorias no contrato de resposta}: Expandir o esquema padronizado do \textit{Orchestrator} para facilitar auditabilidade, rastreamento de decisões e integração com sistemas corporativos.

    \item \textbf{Ciclos contínuos de validação}: Adotar práticas de CI/CD com testes automáticos de segurança, avaliando regressões de desempenho conforme novos modelos ou atualizações são introduzidos.

    \item \textbf{Avaliação em cenários reais}: Implantar a arquitetura em contextos reais — como chatbots institucionais, sistemas educacionais ou plataformas de suporte — para observar comportamento sob demanda real e requisitos operacionais diversos.
\end{itemize}

Assim, este trabalho não apenas demonstra a viabilidade de mecanismos de segurança multicamadas para LLMs, mas também abre caminho para pesquisas futuras que possam ampliar a precisão, reduzir a latência, fortalecer a detecção de vieses e consolidar práticas de IA responsável em ambientes reais.



\bibliographystyle{sbc}
\bibliography{sbc-template}

\end{document}
